% Appendix: implementation details (all hyperparameter values; symbols as in the method section).
% Values: configs/s0.yaml, configs/s3.yaml, configs/ta.yaml, treeground/{tree,tree_attn,selector,sampling,lora}.py.
\section{Implementation details}
\label{sec:impl}

\paragraph{Candidate pool.}
The point grid has $G=32$ points per side, and we keep $N_g=256$ query-free masks after removing near-duplicates at
$\delta_g=0.9$. For each phrase we add the $k=3$ highest-scoring SAM\,3 instances, with $\delta_t=0.95$. The
containment tree uses $\rho=0.8$. This gives $N\approx260$ candidates per query.

\paragraph{Mask tokens.}
The geometry vector has $d_g=9$ entries (box, log area, log depth, log numbers of children and siblings, log area
ratio to the parent) and the source vector $d_q=k+2=5$ (a flag, a one-hot rank, the score). The MLP $f$ first maps
$[g_v;q_v]$ to $256$ dimensions with one hidden layer, then maps its concatenation with $a_v$ to the width of the
MLLM with one hidden layer, both with GELU.

\paragraph{Attention and selection head.}
The region margin is $d=1$ cell. The tree bias has seven relations (self, parent, child, sibling, ancestor,
descendant, other) and one scalar per relation and layer. $\psi_q$ and $\psi_k$ are two-layer MLPs with hidden
width $1024$ and output width $256$. The scale $\lambda$ is initialised to $1/0.07$ and capped at $100$; the offset
$\beta$ is initialised to $-\log N_g$. We train heads at layers $\ell\in\{12,24,36\}$ of the 36-layer decoder jointly
and choose $\ell^{\star}$ with the checkpoint on validation, which gives $\ell^{\star}=36$.

\paragraph{Training.}
LoRA has rank $16$ ($\alpha=32$, dropout $0.05$) on the attention and MLP projections of the language model. A
candidate is a positive or ignored at $\theta=0.5$. Each query uses $n_h=32$ hard and $n_u=32$ random negatives from
its image and $n_c=64$ candidates from other images. The focal loss has focusing parameter $2$; the family loss uses
$\mu=0.1$ and $w_{\mathrm{fam}}=1$. Training data are PACO-LVIS, Pascal-Part and gRefCOCO (train splits) with
per-step shares $0.6$ / $0.2$ / $0.2$. We train for $12{,}000$ steps with a batch of $32$ images, AdamW, learning
rates $2\times10^{-4}$ (LoRA and the selection-token embedding), $3\times10^{-4}$ (heads and mask-token MLP) and $10^{-2}$ (tree
biases), $500$ warm-up steps followed by cosine decay, and gradient clipping at $1$.

\paragraph{Inference.}
Under the one-mask-per-query protocol the answer is the candidate with the highest score $s_v$. On gRefCOCO the
model answers ``no target'' when $\sigma(s_v)<\tau$ for every candidate, with $\tau=0.925$ chosen on the validation
split. Under the multi-mask protocol (Tab.~\ref{tab:multitarget}) the answer is the union of all candidates with
$\sigma(s_v)\geq t$, with $t=0.9$.

\paragraph{Query rewriting for reasoning queries.}
A reasoning query is a full question, which is not a useful text prompt for SAM\,3. For these queries only, the
base Qwen3-VL-4B (without our LoRA weights) first names the target with the prompt: ``Read the request about this
image and name the one object or object part it refers to. Answer with a short noun phrase only, such as `the red
cup' or `the handle of the door'. Request: \{query\}''. Decoding is greedy with at most $24$ new tokens. The noun
phrase replaces the query's phrases as the SAM\,3 prompt for the text candidates; if the answer is empty we fall
back to the phrases of the original query. The selection model is unchanged and still reads the original question.
This rule was fixed on the ReasonSeg validation split before any test run.
